% AI-written, AI cross-reviewed; no human sign-off; not a claim that RH is proved
% Sections of blueprint v13 with local source-root references normalized. Selected-source hashes identify the original fragments.

% blueprint/ch_O.tex:1-133
\chapter{Overview of the formalization}\label{ch:O}

\section*{Scope and verification}
This blueprint follows five declarations in \path{Main.lean:4--12}, beginning with
\texttt{riemannHypothesis} and ending with \texttt{LFunction\_eq\_zero\_re\_of\_isPrimitive}.
Chapter~\ref{ch:E} gives their precise statements and traces the reductions between them.
The statements, the audit of their dependencies, and the adequacy of the mathematical explanations
in this document are three separate questions. For the first two, Sutherland publishes the
five axiom lists and the results of comparison against a challenge importing Mathlib. The
dated results and the scope of the comparison are given with Theorem~\ref{thm:E-main}
\cite{E-Sutherland}. This is an informal, AI-written draft for human review.
No chapter has been signed off by a human mathematician. The trust ledger on page~\pageref{ch:trust} lists the inputs still awaiting a hand reconstruction.

A proof marked \emph{complete} supplies the argument from its displayed hypotheses and cited
inputs. \emph{Reconstructed} identifies a proof whose main steps were followed but which
delegates intermediate arguments. \emph{Statement only} records a result whose proof is left
to its source. In Chapter~\ref{ch:N}, \emph{unread} identifies predicates left folded in the
758-node catalogue. Named declarations and source locations in the chapters give a starting
point for checking a statement or a delegated step.

\section*{Five layers of the argument}

Alp\"oge describes the number of iterated theta lifts as playing the role of multiplicity in
Stepanov's method. This is the author's interpretation of the construction \cite{O-remark}.
In the formalization, $m$ counts the columns. The chain of $m$ operators, the conditions on
its input, and the inequalities summed over its columns are distinct objects. Following them
from the conclusion back to the hypotheses gives the following route.

\begin{description}[leftmargin=0pt,style=nextline]
\item[E. From windows to a zero-free half-plane.]
The level sum combines a centered principal term with weighted real parts of the nonprincipal
character sums (Definition~\ref{def:E-level}). The orthogonality identity compares a prime sum
with an explicit main term. Given the lattice inequality and the size and character hypotheses,
the choice of $m$ near $\sqrt{x}$ gives the one-scale bound
(Lemma~\ref{lem:E-size}, Proposition~\ref{prop:E-ch-core}). For a family satisfying the window and lattice-chain hypotheses, Proposition~\ref{prop:E-windows-ineq} applies this estimate at every sufficiently large window point.
The analytic step uses windows to control a subsequence of Taylor coefficients. A hypothetical
pole to the right of $\theta$ would force a lower bound at \emph{every sufficiently large order}:
isolate a nearest pole, or use the positive residue in the real-pole case. The two estimates
then conflict on that subsequence. See Theorem~\ref{thm:E-landau-window} and
\texttt{Deduction.landau\_window} (\path{LandauTheorems/landauWindow.lean:83--166}).

\item[L. What is uniform in the end-slope estimates?]
For each estimate on a fixed family one chooses a single constant and threshold before varying the scale and the
member of the family. The mean-slope lower bound still contains $1+m/k$, and the maximal-slope
upper bound for $N_u$ still contains $CmB_0$. In the division estimate,
\texttt{L3main} is one upper bound after summing over all columns.
\texttt{weight\_height\_sum\_le} is an upper bound for each column and place.
Of these three inputs, \texttt{weight\_sum\_eq} is the identity.
For $\log T(u)$ one has a precise lower bound and a separate upper bound $CmL$.
The lower bound enters with a negative coefficient in the assembly. See
Propositions~\ref{prop:L-slopesM} and~\ref{prop:L-slopesN}.
\footnote{\texttt{SlopesMF}, \texttt{SlopesNF}, \texttt{Count\_at}, \texttt{SlopesM\_at},
\texttt{SlopesN\_at}: \path{GlobalSchurBounds/QuadraticLogDerivCongr.lean:167--183,252--254}.
\texttt{DivInputs0R.L3main}: \path{Divisibility/ConstantGeTwentyEight.lean:87}.
\texttt{RecNums.weight\_height\_sum\_le}, \texttt{weight\_sum\_eq}:
\path{Divisibility/Inputs.lean:90--92}.}

\item[G. Where does the stronger local gain occur?]
For an integral input the bare column $K_j$ gives a lower bound involving $g^A$.
At a selected place $v$ the stronger $g^B$ bound applies after the Hecke operator,
to $Z_{v,j+1}K_j$. The correction in both gains can depend on earlier weights through
$w_j-w_{j'}$. The final sum gives a divisibility lower bound on the specified submodule
$\mathrm{Iprod}(S,VV)$. When $VV$ is nonempty, the cost $\tau_P$ contains one
$\mathrm{ep}(p)$ and a separate $m$-fold cost for each selected place. When it is empty,
$\tau_P=0$. See Remark~\ref{rem:G-key}.
\footnote{\texttt{FormalGroupSystem.nu}, \texttt{de}, \texttt{Conclusion2},
\texttt{Conclusion4}, \texttt{LowTermBound}: \path{GrowthBounds/FormalGroupSystem.lean:28--49,649--650,671--679,708--716}.
\texttt{G.tauP}: \path{GrowthBounds/IsResidueClassZero.lean:341--344}.}

\item[S. How do deformation points bound a Selmer length?]
At the Eisenstein point the argument relates a Selmer length to a cotangent length. Given
the model, rigidity and trace hypotheses, a finite set $X$ of admissible representations
supplies the bound
\[
  V_n \ \le\ 2\sum_{\rho\in X}v_\ell\bigl(\operatorname{tr}\rho(\sigma)-t\bigr).
\]
Here $t$ is the specified Eisenstein trace and the representations in $X$ are pairwise
nonisomorphic. The subsequent integer bound makes it possible to sum over primes. The
one-character step uses this finite set and a one-sided inequality; the conjectural picture
of Hecke congruences is only an interpretation. See Proposition~\ref{prop:S-onechar}.
\footnote{\texttt{UpperTheory.V.U3.model\_upper}: \path{UnitaryThreeGroups/ModelUpper.lean:103--111}.
\texttt{UpperTheory.V.Z.J.OneCharU\_s} and \texttt{NumberU3U\_s}:
\path{DenseThresholdCases/NumberU3US.lean:30--59}.}

\item[D. How do the inputs reach the lattice chain?]
The assembly first combines the thresholds on families and then chooses one family satisfying
them. At the Selmer interface the local hypothesis is
$h_{0d}(w)\leq\operatorname{length} H^0_{\rm twist}(w)$. Finiteness of the right side
gives finiteness of the left side, and division by the positive ramification index preserves
the upper bound. Once \texttt{LinkN}, the small-level rank identity and the Selmer sum bound
are available, the pointwise estimate gives \texttt{DIVinN} with constant $|C|+1$.
The construction of the local index estimate is a separate input. See
Lemma~\ref{lem:D-selmer-units} and Propositions~\ref{prop:D-divinn} and~\ref{prop:D-topdeswin}.
\footnote{\texttt{LeTwistInvariantsAtColumns}, \texttt{selmer\_length\_inequality}:
\path{AdmissibleDesigns/HigherColumnInvariantsBound.lean:135,371--405,461--465}.
\texttt{UpperTheory.Top.divinN\_at}, \texttt{divinN\_of\_hypV}:
\path{CharacterBoxes/DivinNOfHypV.lean:76--130}.}
\end{description}

\section*{What the reader can check here}
Chapter~\ref{ch:P} follows the five inputs listed in Chapter~\ref{ch:Z}.
The final column gives the point at which to resume a specialist review.

{\small
\begin{longtable}{@{}p{0.06\linewidth}p{0.43\linewidth}p{0.44\linewidth}@{}}
\toprule
Layer & Argument supplied here & Remaining inputs for review \\\midrule
\endfirsthead
\toprule
Layer & Argument supplied here & Remaining inputs for review \\\midrule
\endhead
E & The pole contradiction and the choice of scale under the lattice
hypothesis (Theorem~\ref{thm:E-landau-window}, Proposition~\ref{prop:E-windows-ineq}). &
Construction of the windows uses the arithmetic and lattice inputs in the
four layers below.\\[4pt]
L & The signed local index, the change of type in the Gram determinant
and the sum giving the dual slope and clause~7 (Z5,
Section~\ref{sec:P-z5}). &
The local budgets, the sum of column costs and the comparison between levels for the specified family. The two height realizations are proved in Section~\ref{sec:P-z5}.\\[4pt]
G & The three jet estimates, the normal form for the biadditive kernel and the integral normalized Hecke action under admission (Sections~\ref{sec:P-r4-z41} and~\ref{sec:P-r4-z41-kernel}). The conditional four-way split and the local estimates with constants $(11,1)$ appear at \ref{par:r4-z43-deep}. &
Admission on one served chain requires the inputs listed on p.~\pageref{r8:Z4.2.12}. The common columns and integral period comparison remain open (p.~\pageref{r6:review:Z4.4}). Z4.3 derives its inert CM coordinate. The normalized theta transformation and selected finite expansion remain open (p.~\pageref{r9:z43-open}).\\[4pt]
S & The graded model and both cotangent length bounds under the shared arithmetic inputs (Section~\ref{sec:P-z14-full}); cycles with coefficients modulo $q$ on integral matrices and the quantified torsion bound (Section~\ref{sec:P-r4-K2}, Proposition~\ref{prop:P-r4-k7}). &
The stable odd-rank evaluations and ordinary tangent argument appear in Sections~\ref{sec:P-r5-k4} and~\ref{sec:P-r5-z12}. The status of ordinary rigidity is partial (p.~\pageref{r6:review:Z1.2}). Dense automorphy and the global Kato class still have open components. The cohomological count is on page~\pageref{r6:Z1.5.02}.\\[4pt]
D & Construction of the normalized integral presentation, the arithmetic row comparison, both public lifts and the off-set dimension count give the eight-term bound from the public hypotheses (Section~\ref{sec:P-r4-Z32}). &
When applying the bound, verify the local arithmetic, coefficient, topology and tame hypotheses at that particular family. Chapter~\ref{ch:Z} identifies the ranges to inspect.\\
\bottomrule
\end{longtable}}

Of the twenty original recursive inputs, \TrustOriginalComplete{} have hand deductions from their stated hypotheses. Across the \TrustUnits{} expert assignments, \TrustComplete{} are complete, \TrustPartial{} are partial and \TrustStructure{} have a structure-only map. The marking and Euler-system hypotheses are compared in Section~\ref{r6:publications}. The nine questions in Section~\ref{r8:kato-questions} specify the remaining Kato comparisons; Section~\ref{r7:z42-common-operator} gives the contraction for the two adjacent site frames. Section~\ref{sec:T-r5-five} gives the five largest project-specific theta arguments and their open steps. The expert tables in Chapter~\ref{ch:Z} list
individual reading assignments. Section~\ref{sec:T-pools} describes the
440 separately assessed declarations belonging to the eleven largest
eligible infrastructure pools in the theta branch.


% blueprint/ch_E.tex:1-236
% Chapter E of the blueprint of L. Alpoge's Lean 4 development (Nov 2026).
% Revision checked against the original pinned Lean source tree and the pinned Mathlib sources.
\chapter{From a Chebyshev-type inequality to GRH and RH}\label{ch:E}

This chapter translates the top of the Lean development of L.~Alp\"oge: the five theorems of
\texttt{Main.lean} and the chain of deductions that leads from them down to the place where one
real inequality between two explicit quantities, written
$\mathrm{ASK}\le \mathrm{MAIN}+\mathrm{error}$, is used. That inequality is a statement about
slopes of Hermitian lattices; it is \emph{not} proved here (it is the subject of Chapter~L, with
inputs from Chapters~S and~D). What is proved here is ordinary analytic number theory:
\begin{enumerate}
\item elementary reductions (functional equation, Euler factors) from one statement about
primitive characters to the five theorems (Section~\ref{sec:E-main});
\item a one-sided theorem of Landau type on \emph{windows}: an upper bound
$\mathcal G(y)\le y^{\theta}$ for a positive combination $\mathcal G$ of smoothed prime sums,
valid only on intervals $[Y,Y^{\lambda}]$ for a sequence $Y\to\infty$, forces all the
$L$-functions involved to be zero-free in $\operatorname{Re}s>\theta$ (Section~\ref{sec:E-landau});
\item an exact orthogonality identity which rewrites $\mathrm{ASK}-\mathrm{MAIN}$ as a character
average of the same smoothed prime sums (Section~\ref{sec:E-balance});
\item a bookkeeping argument (``Chebyshev-type bound'') that converts the inequality
$\mathrm{ASK}\le\mathrm{MAIN}+\mathrm{error}$, summed over $m\approx\sqrt x$ copies, into
$\mathcal G(x)\le x^{1/2+O(\varepsilon)}$ (Sections~\ref{sec:E-core}--\ref{sec:E-coregm});
\item the composition of these steps (Section~\ref{sec:E-top}).
\end{enumerate}
In the author's own outline the last step is described as an explicit formula with a variance test
function. The Lean text at this level does not contain an explicit formula: the passage to zeros
is made by item~(2), and the ``$m\approx\sqrt X$ iterated lifts'' of the outline appear in this
chapter only as an integer $m$ with $x^{1/2-\varepsilon}\le m\le x^{1/2+\varepsilon/2}$ that counts
the copies of the inequality. Nothing in this chapter is claimed beyond what the Lean statements
say; the results are due to Alp\"oge and the proofs below are translations of his Lean proofs,
written out so that they can be checked by hand.

\medskip\noindent\textbf{Conventions.} $\Lambda$ is the von Mangoldt function, $\gamma$ is Euler's
constant, $p$ always denotes a prime. For a Dirichlet character $\chi$ modulo $N$ we write
$L(s,\chi)$ for the meromorphic continuation of its $L$-function; $\chi^{*}$ is the primitive
character inducing $\chi$ and $\operatorname{cond}\chi$ its conductor. The trivial character modulo
$1$ is written $\mathbf 1$, so $L(s,\mathbf 1)=\zeta(s)$. Mathlib's assigned value at $\zeta(1)$ is nonzero
(\texttt{riemannZeta\_ne\_zero\_of\_one\_le\_re}, \nolinkurl{Mathlib/NumberTheory/LSeries/Nonvanishing.lean:411}).
For a general principal character the half-plane statement explicitly excludes $s=1$
(\texttt{DirichletLNonvanishingRightOfHalf}, \nolinkurl{Statements.lean:6}). Levels (finite sets of characters) are written
$\mathcal L$ to avoid a clash with $\Lambda$. At each numbered definition or result the
source footnote gives the files and starting lines in the order of the Lean declarations in the
margin. Paths begin at \texttt{LEAN\_SRC}, or at the Mathlib checkout for
\nolinkurl{Mathlib/} paths.

%======================================================================
\section{The main theorems and the reduction to primitive characters}\label{sec:E-main}

The five theorems of \texttt{Main.lean} are all consequences of one statement, called
\texttt{EveryCharacter} in Lean: no Dirichlet $L$-function vanishes in the open half-plane
$\operatorname{Re}s>\tfrac12$, subject to the principal-character exception at $s=1$. This section states the five theorems and proves the reductions.
They use only the functional equation, the non-vanishing on $\operatorname{Re}s\ge1$, and the
comparison of $L(s,\chi)$ with $L(s,\chi^{*})$; all three are in Mathlib \cite{E-mathlib-dircont,E-mathlib-nonvanishing}.

\begin{definition}[The four target statements]\label{def:E-statements}
\lean{DirichletLNonvanishingRightOfHalf, RiemannZetaZerosOnLine, DirichletLZerosOnLine, PrimitiveDirichletLZerosOnLine, Claim.HalfPlane, Claim.Strip, Claim.Zeta, UpperTheory.EveryCharacter}
\footnote{\nolinkurl{Statements.lean:6,3,10,14}, \nolinkurl{Development/AutomorphicArithmetic/CharacterBoxes/HalfPlane.lean:3,6,9}, \nolinkurl{Development/AutomorphicArithmetic/CharacterBoxes/EveryCharacter.lean:13}.}%
\begin{enumerate}
\item[(HP)] \emph{Half-plane statement.} For every $N\ge1$, every Dirichlet character $\chi$ modulo
$N$ and every $s\in\mathbb C$ with $\operatorname{Re}s>\tfrac12$ and ($\chi\neq\chi_0$ or $s\neq1$),
one has $L(s,\chi)\neq0$. (Here $\chi_0$ is the principal character modulo $N$.) In Lean this is
\texttt{DirichletLNonvanishingRightOfHalf}; \texttt{Claim.HalfPlane} and
\texttt{UpperTheory.EveryCharacter} are the same proposition.
\item[(ST)] \emph{Strip statement.} For every $N\ge1$, $\chi$ modulo $N$ and $s$ with
$0<\operatorname{Re}s<1$: $L(s,\chi)=0$ implies $\operatorname{Re}s=\tfrac12$
(\texttt{DirichletLZerosOnLine}, \texttt{Claim.Strip}).
\item[(PR)] The same as (ST) restricted to primitive $\chi$
(\texttt{PrimitiveDirichletLZerosOnLine}).
\item[(Z)] For every $s\in\mathbb C$: if $\zeta(s)=0$, $s\notin\{-2,-4,-6,\dots\}$ and $s\neq1$, then
$\operatorname{Re}s=\tfrac12$ (\texttt{RiemannZetaZerosOnLine}, \texttt{Claim.Zeta}; this is
literally Mathlib's \texttt{RiemannHypothesis}).
\end{enumerate}
\end{definition}

\begin{lemma}[Euler factors do not vanish]\label{lem:E-euler-factor}
\lean{Claim.factor_ne_zero}
\footnote{\nolinkurl{Development/AutomorphicArithmetic/CharacterBoxes/ClaimStrip.lean:7}.}%
Let $\psi$ be a Dirichlet character, $p$ a prime and $\operatorname{Re}s>0$. Then
$1-\psi(p)p^{-s}\neq0$.
\end{lemma}
\begin{proof}
$|\psi(p)p^{-s}|\le p^{-\operatorname{Re}s}<1$.
\textit{[Status: complete]}
\end{proof}

\begin{proposition}[(HP) implies (ST)]\label{prop:E-strip}
\lean{Claim.strip}\uses{def:E-statements,lem:E-euler-factor}
\footnote{\nolinkurl{Development/AutomorphicArithmetic/CharacterBoxes/ClaimStrip.lean:59}.}%
Statement \textup{(HP)} implies statement \textup{(ST)}; hence also \textup{(PR)}.
\end{proposition}
\begin{proof}
Let $\chi$ be a character modulo $N$ and $L(s,\chi)=0$ with $0<\operatorname{Re}s<1$; in particular
$s\neq0,1$.

If $\operatorname{Re}s>\tfrac12$ this contradicts (HP) (the side condition holds because $s\neq1$).

Suppose $\operatorname{Re}s<\tfrac12$. Let $\chi^{*}$ be the primitive character modulo
$f=\operatorname{cond}\chi$ inducing $\chi$. For $s\neq1$,
\begin{equation}\label{eq:E-changelevel}
L(s,\chi)=L(s,\chi^{*})\prod_{p\mid N}\bigl(1-\chi^{*}(p)p^{-s}\bigr)
\end{equation}
(\texttt{DirichletCharacter.LFunction\_changeLevel} in \cite{E-mathlib-dircont}). By
Lemma~\ref{lem:E-euler-factor} the product is non-zero, so $L(s,\chi^{*})=0$
(\texttt{Claim.prim\_zero}). Let $\Lambda(s,\psi)=L(s,\psi)\,\gamma_\psi(s)$ be the completed
$L$-function of a primitive $\psi$, where $\gamma_\psi(s)=\pi^{-s/2}\Gamma(s/2)$ for even $\psi$ and
$\pi^{-(s+1)/2}\Gamma((s+1)/2)$ for odd $\psi$. The factor $\gamma_\psi$ is finite and non-zero for
$\operatorname{Re}s>0$, hence $\Lambda(s,\chi^{*})=0$. The functional equation
\[
\Lambda(1-s,\psi^{-1})=f^{\,s-1/2}\,\varepsilon(\psi^{-1})\,\Lambda(s,\psi)\qquad(\psi\text{ primitive modulo }f)
\]
(\texttt{DirichletCharacter.IsPrimitive.completedLFunction\_one\_sub} in \cite{E-mathlib-dircont},
applied to $\psi^{-1}$) gives $\Lambda(1-s,(\chi^{*})^{-1})=0$, and since
$\operatorname{Re}(1-s)>0$ the gamma factor is again finite and non-zero, so
$L(1-s,(\chi^{*})^{-1})=0$ (\texttt{Claim.prim\_reflect}). But
$\operatorname{Re}(1-s)>\tfrac12$ and $1-s\neq1$, which contradicts (HP) for the character
$(\chi^{*})^{-1}$ modulo $f$. Hence $\operatorname{Re}s=\tfrac12$.
\textit{[Status: complete]}
\end{proof}

\begin{proposition}[(HP) implies the Riemann Hypothesis]\label{prop:E-zeta}
\lean{Claim.zeta}\uses{prop:E-strip}
\footnote{\nolinkurl{Development/AutomorphicArithmetic/CharacterBoxes/ClaimStrip.lean:75}.}%
Statement \textup{(HP)} implies statement \textup{(Z)}.
\end{proposition}
\begin{proof}
Let $\zeta(s)=0$ with $s\neq1$ and $s$ not a negative even integer. Since $\zeta(0)=-\tfrac12$,
$s\neq0$.

\emph{Case $0<\operatorname{Re}s<1$.} As $\zeta(s)=L(s,\mathbf 1)$, Proposition~\ref{prop:E-strip}
with $N=1$ gives $\operatorname{Re}s=\tfrac12$.

\emph{Case $\operatorname{Re}s\ge1$.} Impossible, because $\zeta$ has no zero there
(\texttt{riemannZeta\_ne\_zero\_of\_one\_le\_re}, \cite{E-mathlib-nonvanishing}).

\emph{Case $\operatorname{Re}s\le0$.} Put $\xi(s)=\pi^{-s/2}\Gamma(s/2)\zeta(s)$. The factor
$\pi^{-s/2}\Gamma(s/2)$ is finite and non-zero unless $s/2\in\{0,-1,-2,\dots\}$, which is excluded;
hence $\xi(s)=0$. By the functional equation $\xi(1-s)=\xi(s)$ we get $\xi(1-s)=0$, and
$\operatorname{Re}(1-s)\ge1$, where $\pi^{-(1-s)/2}\Gamma((1-s)/2)$ is finite and non-zero and
$1-s\neq1$. Hence $\zeta(1-s)=0$ with $\operatorname{Re}(1-s)\ge1$, which is impossible as in the
previous case.
\textit{[Status: complete]}
\end{proof}

The Lean development does not prove (HP) directly; it proves the following statement about
primitive characters and deduces (HP) from it.

\begin{definition}[Primitive characters, zero-freeness]\label{def:E-pchar}
\lean{Formalism.HC.PChar, Formalism.HC.PChar.of, Formalism.HC.PChar.ZeroFreeAbove, Formalism.HC.ZeroFreeAll}
\footnote{\nolinkurl{Formalism/SharpRestriction/RamificationFrobenius/TotientGaloisClasses/PrimePairSizeCondition.lean:7,32}, \nolinkurl{Formalism/SharpRestriction/RamificationFrobenius/TotientGaloisClasses/ExplicitFormulaErrorOfHostConditions.lean:11,18}.}%
A \emph{primitive character} (\texttt{PChar}) is a pair $\psi=(N,\chi)$ of an integer $N\ge1$ and a
primitive Dirichlet character $\chi$ modulo $N$; we simply write $\psi$ for $\chi$. The set of all
primitive characters is denoted $\mathcal P$; it contains $\mathbf 1=(1,1)$. For an arbitrary
character $\chi$ modulo $N$ we write $\chi^{*}\in\mathcal P$ for the primitive character inducing it
(\texttt{PChar.of}). For $\theta\in\mathbb R$, $\psi$ is \emph{zero-free above $\theta$} if
$L(s,\psi)\neq0$ for all $s$ with $\operatorname{Re}s>\theta$ (for $\psi=\mathbf 1$ the point $s=1$
is not a zero, see the conventions). The statement \texttt{ZeroFreeAll} is: every
$\psi\in\mathcal P$ is zero-free above $\tfrac12$.
\end{definition}

\begin{proposition}[From primitive to arbitrary characters]\label{prop:E-imprimitive}
\lean{UpperTheory.Top.imprimitive}\uses{def:E-pchar,def:E-statements,lem:E-euler-factor}
\footnote{\nolinkurl{Development/AutomorphicArithmetic/CharacterBoxes/Imprimitive.lean:5}.}%
\texttt{ZeroFreeAll} implies \textup{(HP)}.
\end{proposition}
\begin{proof}
Let $\chi$ be a character modulo $N$, $\operatorname{Re}s>\tfrac12$, and $\chi\neq\chi_0$ or
$s\neq1$. If $\chi\neq\chi_0$ then $\chi^{*}\neq\mathbf 1$ (a character induced by $\mathbf 1$ is
principal). So $\chi^{*}\neq\mathbf1$ or $s\neq1$, and \eqref{eq:E-changelevel} holds at $s$. The
factor $L(s,\chi^{*})$ is non-zero by hypothesis and the finite product is non-zero by
Lemma~\ref{lem:E-euler-factor}.
\textit{[Status: complete]}
\end{proof}

\begin{theorem}[\texttt{everyCharacter}; Alp\"oge, formalised]\label{thm:E-every-character}
\lean{UpperTheory.FinalTheorems.everyCharacter}\uses{prop:E-imprimitive,thm:E-zero-free-all}
\footnote{\nolinkurl{Development/AutomorphicArithmetic/CharacterBoxes/MainConclusion.lean:10}.}%
Statement \textup{(HP)} holds.
\end{theorem}
\begin{proof}
The published wrapper \texttt{UpperTheory.FinalTheorems.everyCharacter} invokes
\texttt{everyCharacter\_of\_row}, which calls
\texttt{every\_character\_of\_literature\_of\_vWin}
(\nolinkurl{Development/AutomorphicArithmetic/CharacterBoxes/MainConclusion.lean:7--11}).
The analytic part of this assembly uses
\texttt{UpperTheory.Top.every\_character\_of\_green\_v\_win}, which is
Proposition~\ref{prop:E-imprimitive} applied to
\texttt{UpperTheory.Top.zero\_free\_all\_of\_green\_v\_win}. The latter is
Theorem~\ref{thm:E-zero-free-all} below with $\lambda=\tfrac32$, applied to the two inputs
\textup{(I1)} and \textup{(I2)} of Section~\ref{sec:E-top}: the statement \texttt{TopG}
(Definition~\ref{def:E-topG}) and, for every $\psi\in\mathcal P$, the statement
$\texttt{HostExW}(\tfrac32,\texttt{GChi}(\psi))$ (Definition~\ref{def:E-hostexw}). These two inputs
are theorems of the Lean development proved outside this chapter (lattice chain: Chapter~L; Selmer
groups: Chapter~S; design and interface: Chapter~D); see Remark~\ref{rem:E-inputs} for the Lean
declarations that supply them.
\textit{[Status: complete relative to the inputs \textup{(I1)}, \textup{(I2)}, which are proved in
Chapters L, S and D]}
\end{proof}

\begin{theorem}[The five theorems of \texttt{Main.lean}]\label{thm:E-main}
\lean{LFunction_ne_zero_of_half_lt_re, riemannHypothesis, riemannZeta_eq_zero_re, LFunction_eq_zero_re, LFunction_eq_zero_re_of_isPrimitive}
\uses{thm:E-every-character,prop:E-strip,prop:E-zeta}
\footnote{\nolinkurl{Main.lean:8,4,6,10,12}.}%
The five declarations in \texttt{Main.lean} have the following conclusions.
\begin{enumerate}
\item \texttt{LFunction\_ne\_zero\_of\_half\_lt\_re}: statement \textup{(HP)};
\item \texttt{riemannHypothesis}: Mathlib's \texttt{RiemannHypothesis}, i.e.\ \textup{(Z)};
\item \texttt{riemannZeta\_eq\_zero\_re}: \textup{(Z)} in the unfolded form of
Definition~\ref{def:E-statements};
\item \texttt{LFunction\_eq\_zero\_re}: statement \textup{(ST)};
\item \texttt{LFunction\_eq\_zero\_re\_of\_isPrimitive}: statement \textup{(PR)}.
\end{enumerate}
\end{theorem}
\begin{proof}
For (1), \texttt{allChi} supplies Theorem~\ref{thm:E-every-character}. For (2) and (3),
\texttt{rhOfAll} applies (HP) to the character of modulus one and invokes the zeta reduction.
For (4), the second component of \texttt{claimAll} supplies \texttt{Claim.strip}. For (5),
\texttt{critAll} uses primitivity to reflect a zero on the left via the functional equation for the
inverse character, and applies (HP) on the right.
\phantomsection\label{r10:main-wrappers}The wrappers are \texttt{allChi}, \texttt{rhOfAll}, \texttt{critAll}, \texttt{claimAll}
(\nolinkurl{Supplements/Kernels/MainTheorems.lean:7,9,11,15}). The two direct reductions are
\texttt{Formalism.VW.rh\_of\_every} and \texttt{Formalism.VW.critLine}
(\nolinkurl{Formalism/GammaFactorNonvanishing/CritLine.lean:29,11}).
\textit{[Status: complete relative to Theorem~\ref{thm:E-every-character}]}
\end{proof}

\paragraph{External axiom and statement checks.}
\input{r10_audit}
Sutherland's axiom log lists \texttt{propext}, \texttt{Classical.choice} and \texttt{Quot.sound}
for each of the five declarations. The \texttt{con-leche --compare} log gives the
individual matches to the challenge. The standalone Comparator v4.33.0 entry uses
the original build and a local compilation of \texttt{Challenge.lean}
\cite{E-Sutherland,E-audit,E-compare}.
The local transcript \nolinkurl{out/S1_main_axioms.txt:1} stops at a missing object file before the
axiom lists are printed.

%======================================================================
